\section*{Реферат}
\addcontentsline{toc}{section}{Реферат}

Выпускная квалификационная работа посвящена разработке Android-приложения для совместного планирования поездок. Актуальность исследования обусловлена тем, что в реальной практике групповой подготовки поездки пользователи распределяют задачи между несколькими сервисами, что приводит к дублированию данных и потере контекста решений.

Объектом исследования является процесс совместного планирования поездок в малых группах пользователей. Предметом исследования выступают пользовательские сценарии, функциональные требования и проектные решения мобильного приложения, объединяющего ключевые этапы планирования в едином цифровом пространстве.

Цель работы состоит в проектировании и реализации мобильного приложения, обеспечивающего целостный сценарий группового планирования поездки. В рамках работы проведен анализ предметной области и существующих решений, сформулирован функциональный разрыв, разработаны архитектурные и модульные решения, а также описана программная реализация подсистем приложения.

Результатом работы является структурированная модель приложения, включающая управление поездками и участниками, работу с идеями и маршрутом, учет расходов и балансов, погодный контекст, AI-предложения, настройки уведомлений, а также офлайн-режим с последующей синхронизацией данных.

Практическая значимость работы заключается в применимости полученного результата для повседневного группового планирования поездок и в возможности дальнейшего развития системы в рабочей среде.

\textbf{Ключевые слова:} Android-приложение, совместное планирование поездок, маршрут по дням, учет расходов, баланс участников, офлайн-синхронизация, AI-предложения.
